% ============================================================
% Capítulo 9 — Losses: conceptos + código
% ============================================================
\chapter{Funciones de pérdida: triplet y cross-entropy}
\label{cap:losses}

Este capítulo explica los conceptos detrás de las dos pérdidas
usadas y muestra su código línea por línea.

\section{¿Qué es un embedding y por qué necesita pérdidas?}

Después del forward tenemos dos vectores por muestra:
\begin{itemize}
  \item $\mathbf{I}' \in \mathbb{R}^{1024}$ (imagen).
  \item $\mathbf{T}' \in \mathbb{R}^{1024}$ (texto).
\end{itemize}

\noindent\textbf{Objetivo}: vectores de la misma persona deben
estar \emph{cerca}, y de personas diferentes \emph{lejos}.

\noindent\textbf{Problema}: el modelo no sabe cómo hacerlo solo.
Necesitamos una función que mida qué tan mal lo está haciendo y
guíe el aprendizaje: la \textbf{pérdida} (loss).

\section{¿Qué es la triplet loss?}

\subsection{Concepto}

La \emph{triplet loss}~\cite{schroff2015facenet} trabaja con
\emph{triplets}: tríos de muestras (anchor, positive, negative).

\begin{itemize}
  \item \textbf{Anchor} (A): muestra de referencia.
  \item \textbf{Positive} (P): muestra del \emph{mismo} pid.
  \item \textbf{Negative} (N): muestra de \emph{otro} pid.
\end{itemize}

La pérdida enforces que:

\begin{equation}
  \text{sim}(A, P) - \text{sim}(A, N) + \text{margin} \leq 0
  \label{eq:triplet}
\end{equation}

\noindent\textbf{En palabras}: el anchor debe estar \emph{más
cerca} del positive que del negative, con un margen mínimo
($\text{margin} = 0.5$ en este proyecto).

\subsection{¿Por qué el margen?}

Sin margen, basta con que $\text{sim}(A, P) > \text{sim}(A, N)$.
Esto da una solución trivial: todos los embeddings en un punto.
Con margen exigimos que la diferencia sea $\geq \text{margin}$,
lo que \emph{obliga} a aprender embeddings informativos.

\section{¿Qué es semi-hard mining?}

\subsection{Los tres tipos de negativos}

Para un anchor dado, hay tres tipos de negativos:

\begin{enumerate}
  \item \textbf{Easy}: $\text{sim}(A, N) \ll \text{sim}(A, P)$.
        El modelo ya los clasifica bien. \emph{No ayudan} al
        aprendizaje.
  \item \textbf{Hard}: $\text{sim}(A, N) > \text{sim}(A, P)$.
        El modelo ya los confunde. Contribuyen mucho al
        gradiente pero pueden hacer el entrenamiento inestable.
  \item \textbf{Semi-hard}: $\text{sim}(A, P) < \text{sim}(A, N)
        < \text{sim}(A, P) + \text{margin}$. El modelo está
        cerca pero no del todo mal. \emph{Ideal}: gradiente
        útil + entrenamiento estable.
\end{enumerate}

\subsection{Criterio usado en este proyecto}

El código implementa semi-hard mining~\cite{schroff2015facenet}:
\begin{itemize}
  \item Si hay negativos semi-hard disponibles, elegir uno
    aleatoriamente.
  \item Si no hay ninguno, usar uno cualquiera (caso
    degenerado).
\end{itemize}

\noindent\textbf{Ventaja}: estable y eficiente; converge más
rápido que easy mining sin los problemas de hard mining.

\section{¿Qué es la cross-entropy?}

\subsection{Clasificación como identidad}

Otra forma de entrenar el modelo es agregar una cabeza de
\textbf{clasificación} que predice el pid (identidad de persona).
Si el modelo predice correctamente, está discriminando bien.

La cross-entropy mide qué tan lejos están las predicciones del
label correcto:

\begin{equation}
  \mathcal{L}_{\text{CE}} = -\sum_{c=1}^{C} y_c \cdot \log(\hat{y}_c)
  \label{eq:ce}
\end{equation}

\noindent donde $y_c$ es 1 si $c$ es el pid correcto, $\hat{y}_c$ es
la probabilidad predicha para la clase $c$.

\subsection{PyTorch lo hace solo}

\texttt{nn.CrossEntropyLoss(reduction='mean')} computa la CE
promediada sobre el batch. Internamente aplica
\texttt{LogSoftmax + NLLLoss}.

\section{¿Por qué simétrica (i2t + t2i)?}

Tanto RankingLoss como IdentityLoss se calculan en
\textbf{ambas direcciones}:
\begin{itemize}
  \item \textbf{i2t}: predice la clase del texto desde el
    embedding de la imagen.
  \item \textbf{t2i}: predice la clase de la imagen desde el
    embedding del texto.
\end{itemize}

\noindent\textbf{¿Por qué?} El modelo debe ser bueno
\emph{simétricamente}: si una imagen $\leftrightarrow$ texto
describe a la misma persona, ambos sentidos deben predecir el
mismo pid. Esto fuerza coherencia en ambos lados del embedding.

\section{Pérdida total}

\begin{equation}
  \mathcal{L}_{\text{total}} = \alpha \cdot \mathcal{L}_{\text{rank}}
                             + \beta  \cdot \mathcal{L}_{\text{ID}}
\end{equation}

\noindent En este proyecto $\alpha = \beta = 1.0$
(\fline{config.py}{87-88}).

\section{Código: \texttt{calculate\_similarity}}

\begin{minted}[language=Python, numbers, firstnumber=11, linerange=11-18]{python}
def calculate_similarity(image_embedding, text_embedding):
    """Calcula la matriz de similitud coseno entre imágenes y textos."""

    # Aplanar embeddings a (B, 1024) por si vienen con más dims
    image_embedding = image_embedding.view(image_embedding.size(0), -1)
    text_embedding = text_embedding.view(text_embedding.size(0), -1)

    # Normalizar L2 (norma = 1)
    # +1e-8 para evitar división por cero
    image_embedding_norm = image_embedding / (
        image_embedding.norm(dim=1, keepdim=True) + 1e-8
    )
    text_embedding_norm = text_embedding / (
        text_embedding.norm(dim=1, keepdim=True) + 1e-8
    )

    # Producto matricial: similitud entre cada par (img_i, txt_j)
    similarity = torch.mm(image_embedding_norm, text_embedding_norm.t())
    return similarity
\end{minted}

\section{Código: \texttt{RankingLoss}}

\begin{minted}[language=Python, numbers, firstnumber=21, linerange=21-51]{python}
class RankingLoss(nn.Module):
    """Triplet loss con minería semi-hard."""

    def __init__(self, config=None):
        super(RankingLoss, self).__init__()
        self.config = config

    def semi_hard_negative(self, loss):
        """Elige un negativo semi-hard al azar."""
        # loss: array con losses para cada negativo
        # logical_and(0 < loss < margin)
        negative_index = np.where(
            np.logical_and(loss < self.config.margin, loss > 0)
        )[0]
        return np.random.choice(negative_index) if len(negative_index) > 0 else None

    def get_triplets(self, similarity, labels):
        """Para cada anchor, busca un par (positivo, negativo)."""
        similarity = similarity.cpu().data.numpy()
        labels = labels.cpu().data.numpy()
        triplets = []

        for idx, label in enumerate(labels):
            # Índices con OTRO pid (negativos)
            negative = np.where(labels != label)[0]
            # Similitud del anchor consigo mismo (~1)
            ap_sim = similarity[idx, idx]
            # loss = sim(a,n) - sim(a,p) + margin
            loss = similarity[idx, negative] - ap_sim + self.config.margin
            # Elegir un negativo semi-hard
            negetive_index = self.semi_hard_negative(loss)
            if negetive_index is not None:
                triplets.append([idx, idx, negative[negetive_index]])

        if len(triplets) == 0:
            triplets.append([idx, idx, negative[0]])
        triplets = np.array(triplets)
        return torch.LongTensor(triplets)
\end{minted}

\begin{minted}[language=Python, numbers, firstnumber=54, linerange=54-72]{python}
    def forward(self, img, txt, label, sim_neg=None):
        """img: (B,1024), txt: (B,1024), label: (B,)"""
        similarity = calculate_similarity(img, txt)
        if sim_neg is not None:
            similarity = similarity + sim_neg

        # Triplets image->text
        image_triplets = self.get_triplets(similarity, label)
        image_anchor_loss = F.relu(
            self.config.margin
            - similarity[image_triplets[:, 0], image_triplets[:, 1]]
            + similarity[image_triplets[:, 0], image_triplets[:, 2]]
        )

        # Triplets text->image (transponer)
        text_triplets = self.get_triplets(similarity.t(), label)
        similarity = similarity.t()
        texy_anchor_loss = F.relu(
            self.config.margin
            - similarity[text_triplets[:, 0], text_triplets[:, 1]]
            + similarity[text_triplets[:, 0], text_triplets[:, 2]]
        )

        loss = torch.sum(image_anchor_loss) + torch.sum(texy_anchor_loss)
        return loss
\end{minted}

\section{Código: \texttt{IdentityLoss}}

\begin{minted}[language=Python, numbers, firstnumber=10, linerange=10-27]{python}
def weights_init_classifier(m):
    """Inicializa capas Linear con N(0, 0.001) y bias=0."""
    classname = m.__class__.__name__
    if classname.find('Linear') != -1:
        init.normal_(m.weight.data, std=0.001)
        init.constant_(m.bias.data, 0.0)


class classifier(nn.Module):
    """Capa lineal simple: input_dim -> output_dim."""

    def __init__(self, input_dim, output_dim):
        super(classifier, self).__init__()
        self.block = nn.Linear(input_dim, output_dim)
        self.block.apply(weights_init_classifier)

    def forward(self, x):
        return self.block(x)
\end{minted}

\begin{minted}[language=Python, numbers, firstnumber=30, linerange=30-56]{python}
class IdentityLoss(nn.Module):
    """Cross-entropy simétrica: image->text + text->image."""

    def __init__(self, config=None, part=1, class_num=None):
        super(IdentityLoss, self).__init__()
        self.part = part

        # Crear `part` clasificadores (en este proyecto part=1)
        W = []
        for i in range(part):
            W.append(classifier(config.feature_length, class_num))
        self.W = nn.Sequential(*W)

    def calculate_IdLoss(self, image_embedding_local,
                         text_embedding_local, label):
        """image_embedding_local: (B,1024,1), text_embedding_local: (B,1024,1)"""
        label = label.view(label.size(0))
        criterion = nn.CrossEntropyLoss(reduction='mean')

        Lipt_local = 0  # image->text loss
        Ltpi_local = 0  # text->image loss

        for i in range(self.part):
            # Proyectar embedding al espacio de clases
            score_i2t_local_i = self.W[i](image_embedding_local[:, :, i])
            score_t2i_local_i = self.W[i](text_embedding_local[:, :, i])

            Lipt_local += criterion(score_i2t_local_i, label)
            Ltpi_local += criterion(score_t2i_local_i, label)

        loss = (Lipt_local + Ltpi_local) / self.part
        return loss

    def forward(self, image_embedding_local,
                text_embedding_local, label):
        return self.calculate_IdLoss(image_embedding_local,
                                     text_embedding_local, label)
\end{minted}

\section{Uso conjunto en \texttt{train.py}}

\begin{minted}[language=Python, numbers, firstnumber=130, linerange=130-150]{python}
with torch.autocast(device_type=config.device, dtype=torch.float16):
    # Forward
    visual_embeddings, textual_embeddings = model(
        image=preprocessed_images,
        text_ids=token_ids,
        text_length=orig_token_length
    )

    # Aplanar a (B, 1024) - RankingLoss espera aplanados
    visual_emb_flat = visual_embeddings.view(visual_embeddings.size(0), -1)
    text_emb_flat = textual_embeddings.view(textual_embeddings.size(0), -1)
    labels_flat = labels.view(-1)

    # RankingLoss: triplet semi-hard
    ranking_loss = ranking_loss_fnx(visual_emb_flat, text_emb_flat, labels_flat)

    # IdentityLoss: cross-entropy simétrica
    identity_loss = identity_loss_fnx(visual_embeddings, textual_embeddings, labels_flat)

    # Pérdida total: combinación ponderada
    total_loss = (config.ranking_loss_alpha * ranking_loss
                  + config.identity_loss_beta * identity_loss)
\end{minted}
